\documentclass[11pt]{report}
\input{../pri-end/T0_preamble_standalone_A4_De}
% T0_preamble_patches.tex — LOKALER PLATZHALTER
% Im Repo durch die offizielle Version ersetzen.


\title{Dok.~314 -- Das Gitter im Hilbertraum:\\
Geometrie, Spektrum und Faser}
\author{Johann Pascher}
\date{August 2026}

\begin{document}
\maketitle
\tableofcontents

% ============================================================
\chapter*{Dok.~314 -- Das Gitter im Hilbertraum}
\addcontentsline{toc}{chapter}{Dok.~314 -- Das Gitter im Hilbertraum}
% ============================================================

\section*{Worum es geht}

Die FFGFT rechnet geometrisch auf T$^4$ mit dem D4-Gitter und
quantenmechanisch auf $\mathcal{H} = L^2(\text{T}^4)\otimes
\mathbb{C}^3$ (Dok.~230/282). Beim Übergang stellt sich eine
naheliegende Frage: \emph{Werden die Gitter dabei flach?} Bleibt von
Kusszahl~24, Packungsdichte und Trialität etwas übrig, oder ist im
Hilbertraum jedes Gitter dasselbe?

Die Antwort lautet: Der Torus ist ohnehin flach -- die Krümmung war
nie der Unterschied --, und die Gitterstruktur geht nicht verloren,
sondern \textbf{wandert ins Spektrum und in die Faser}. Dieses
Dokument rechnet das durch (drei Skripte, Kap.~J) und findet:

\begin{enumerate}
  \item \textbf{Das Spektrum trennt, was die Krümmung nicht trennt
        [K]:} Z$^4$ beginnt bei $|k|^2 = 1$ mit Entartung~8, D4 bei
        $|k|^2 = 2$ mit Entartung~24; die ungeraden Normen fehlen
        bei D4 vollständig. Die erste Anregung trägt exakt
        $24 = 12+12$ -- FCC-Hälfte und Wicklungshälfte getrennt
        adressierbar.
  \item \textbf{Die Trialität ist eingebaut, nicht angeheftet [K]:}
        $|\text{Aut}(D4)| = 1152 = |W(F_4)|$, Quotient zu $W(D4)$
        ist $6 = |S_3|$; unter den 80 Ordnung-3-Elementen erzeugen
        16 exakt den T$^4/Z_3$-Orbifold mit $|\det(1-A)| = 9$
        Fixpunkten und wirken auf jedem Modenorbit als
        $\mathbb{C}^3$-Zirkulant.
  \item \textbf{Ein Negativbefund mit Satzcharakter [K]:} $5$ teilt
        $1152 = 2^7\cdot3^2$ nicht, also enthält $\text{Aut}(D4)$
        \emph{kein} Element der Ordnung~5. Die Fünffachdrehung, aus
        der $\theta = 2/9$ stammt (Dok.~293), ist mit dem Gitter
        prinzipiell unverträglich -- die Phase gehört einer
        anderen Schicht.
  \item \textbf{Die Verdopplung ist dennoch vorbereitet [K]:}
        $24 = 8\times3 = 4\times6$; vier disjunkte Sechserblöcke mit
        genau dem $C_3$-Inhalt, den $3\oplus3'$ verlangt.
  \item \textbf{Die Reziprozität ist modenweise exakt -- als
        \emph{Ort}-Relation [K]:} $\lambda_4\cdot m = 2\pi$
        identisch für jede Mode (de Broglie auf der
        Massenrichtung; Wicklungszahl $=$ Massenzahl). Die
        \emph{zeitliche} Ablesung trägt $1/\gamma$, das
        Massengemisch einen Jensen-Exzess $\ge1$: exakt wird die
        Relation nur für die ruhende, scharfe Mode.
  \item \textbf{Die Störungsrechnung ist ausgeführt [K]:}
        Multiplettgrößen folgen den Irrep-Dimensionen der
        respektierten Symmetrie, nicht den Orbits --
        $9+8+4+2+1$ im voll symmetrischen Fall, $8\times1+8\times2$
        bei $Z_3$, generisch 24~Einzelniveaus.
  \item \textbf{Die Schließungsgabelung ist spektral unsichtbar
        [K]:} Im eingerollten Zustand gibt es keine fraktalen
        Korrekturen -- Wicklungszahlen sind topologisch, der
        Defekt $d = 100\xi$ akkumuliert nur ausgerollt. Die
        Fälle A/B/C aus Dok.~313/295 lassen das Modenspektrum
        unberührt; die Kapitel~E--H sind damit fallunabhängig.
        Dieselbe Grenze regelt die fraktale Frage (Kap.~J):
        $K_\text{frak}$ ist kumulativ (Dok.~133) und greift nur
        über die Skalenbrücke; lokal wirkt allein
        $D_f^\text{Raum} = 3-\xi$ mit $6{,}7\times10^{-5}$.
\end{enumerate}

Kein Teil der A-Serie; Status: Arbeitsdokument. Bezüge: Dok.~230/282
($\mathcal{H} = L^2(\text{T}^4)\otimes\mathbb{C}^3$), 285
(Verdopplung $6 = 3\oplus3'$), 291 (Kanalanalyse), 293
($\theta = 2/9$ ikosaedrisch), 133 (Herleitung $K_\text{frak}$: kumulativ, RG-Lauf),
295/313 (Schließungsgabelung), 306 (native
Zeit-Energie-Reziprozität), 311 (Vier auf Drei, $24 = 12+12$).

% ============================================================
\section*{Leseführung: der rote Faden}
% ============================================================

\textbf{Die Frage.} Im Hilbertraum wird das Gitter zur Indexmenge
der Fourier-Moden. Damit ist es keine Kugelpackung mehr -- in diesem
Sinn wird jedes Gitter \enquote{flach}. Die Frage ist, ob damit auch
die Information verschwindet.

\textbf{Die Antwort in einer Zeile.} \emph{Geometrie $\to$ Spektrum
$+$ Faser.} Die Entartungen des Laplace-Operators sind die
Thetareihe des Gitters; die Kusszahl wird zur Entartung der ersten
Anregung; die Trialität wird zur Orbifold- und Faserstruktur.

\textbf{Was daraus folgt.} Erstens sind zwei flache Tori am
Spektrum unterscheidbar (Kap.~B). Zweitens wird die
$\tilde T\cdot m = 1$-Struktur bei $R_4 \ne R_{123}$ im Spektrum
ablesbar, und die Niveauordnung ändert sich an genau zwei
Wurzelradien (Kap.~C). Drittens braucht die $Z_3$ nicht gesetzt zu
werden -- sie ist Gittersymmetrie (Kap.~E).

\textbf{Die Grenze, sauber gezogen.} Die Fünffachsymmetrie ist am
Gitter unmöglich (Kap.~F). Das trifft die Verdopplung nicht, weil
diese auf der \emph{Faser} lebt und die Schnittstelle beider Ebenen
genau $C_3$ ist -- besetzt durch die Trialität (Kap.~G). Und die
Container-Struktur gilt auf \emph{jeder} Schale, kann also keine
auszeichnen (Kap.~H) -- dieselbe Muster-gegen-Besetzung-Logik wie in
Dok.~293.

\textbf{Der Vorbehalt, präzisiert.} Gerechnet wird mit dem glatten
Laplace. Kap.~J trennt, was das kostet, und stützt sich dabei auf
Dok.~133: $D_f^\text{Raum} = 3-\xi$ ist \emph{lokal} und wirkt mit
$6{,}7\times10^{-5}$; $K_\text{frak} = 1-100\xi$ ist
\emph{kumulativ} (RG-Lauf über 17 Größenordnungen) und greift
allein über die Skalenbrücke. Nackt sind damit nur die
SI-Absolutwerte ($-1{,}33\,\%$ über den Anker); Zählungen,
Verhältnisse und Container sind es nicht. Die Entartungsaufspaltung
unter Störung ist gerechnet -- sie folgt den Irrep-Dimensionen der
Störungssymmetrie, nicht den Orbits, und liegt bei $\sim10^{-5}$. Kap.~D2 trägt zudem die Grenze zwischen eingerolltem und
ausgerolltem Zustand nach: fraktale Korrekturen gibt es nur
ausgerollt, weshalb die Schließungsgabelung (Dok.~313/295)
spektral unsichtbar bleibt und die Kap.~E--H nicht belastet.

\medskip
\noindent\textbf{In einem Satz:} Im Hilbertraum wird das Gitter
flach, aber nicht folgenlos -- was die Geometrie trug, tragen
Spektrum und Faser, und die Arbeitsteilung zwischen beiden ist
exakt bestimmbar, solange man die glatte Näherung als solche
buchführt.

% ============================================================
\section*{Zeichenerklärung}
% ============================================================

Die Tabelle erklärt die Symbole dieses Dokuments. Größen, die nur an
einer Stelle auftreten, werden dort eingeführt.

\textbf{Notationshinweis.} In der Gittertheorie ist $\Lambda$ das
übliche Zeichen für ein Gitter (Conway/Sloane). Im FFGFT-Korpus ist
$\Lambda$ jedoch belegt -- als kosmologische Konstante, in
$\Lambda$CDM und insbesondere als Abschlussskala $\Lambda^{*} =
1/R_H^2$ (Dok.~312). Dieses Dokument verwendet deshalb
durchgehend $\Gamma$ für das Gitter und $\Gamma^{*}$ für das duale
Gitter. Ebenso wird $L^{*}$ hier \emph{nicht} verwendet, da es in
Dok.~312 die Zyklenlänge bezeichnet; die Kantenlänge des Torus
heißt hier $R_{123}$ bzw.\ $R_4$.

\begin{center}
\renewcommand{\arraystretch}{1.3}
{\small\setlength\tabcolsep{5pt}
\begin{longtable}{lp{9.6cm}}
\toprule
\textbf{Symbol} & \textbf{Bedeutung} \\
\midrule
\endfirsthead
\toprule
\textbf{Symbol} & \textbf{Bedeutung} \\
\midrule
\endhead
\bottomrule
\endfoot
\multicolumn{2}{l}{\textit{Gitter und Torus}} \\
$\Gamma$ & Gitter in $\mathbb{R}^4$; definiert den Torus
  T$^4 = \mathbb{R}^4/\Gamma$ (statt des in der Gitterliteratur
  üblichen $\Lambda$, siehe Notationshinweis) \\
$\Gamma^{*}$ & duales Gitter; trägt die zulässigen Wellenvektoren
  der Fourier-Moden \\
$D4$ & Wurzelgitter $D_4$: alle ganzzahligen $k$ mit
  $\sum k_i$ gerade; Kusszahl 24 \\
$\text{Z}^4$ & kubisches Gitter $\mathbb{Z}^4$; Kusszahl 8,
  exakt selbstdual \\
$D4^{*}$ & duales D4-Gitter; bis auf Skalierung wieder ein
  D4-Gitter \\
$R_{123}$, $R_4$ & Radien der drei Raumrichtungen bzw.\ der
  vierten Richtung \\
$r$ & Radienverhältnis $r = R_4/R_{123}$ (dimensionslos) \\
\midrule
\multicolumn{2}{l}{\textit{Spektrum}} \\
$k$ & Wellenvektor (Modenindex), $k \in \Gamma^{*}$ \\
$k_4$ & vierte Komponente von $k$; $k_4 = 0$ trennt die
  FCC-Hälfte von der Wicklungshälfte \\
$|k|^2$ & Norm; zugleich Eigenwert des Laplace-Operators auf dem
  isotropen Torus \\
$E(k)$ & Eigenwert bei Anisotropie,
  $E = k_1^2+k_2^2+k_3^2+(k_4/r)^2$ \\
$N(n)$ & Entartung der Schale $|k|^2 = n$; Koeffizient der
  Thetareihe \\
$\theta_{D4}$ & Thetareihe des D4-Gitters,
  $\theta_{D4} = (\theta_3^4+\theta_4^4)/2$ \\
$\theta(t)$ & Wärmekern-Thetafunktion, $\theta(t) = \sum_k
  e^{-t|k|^2}$ \\
$\zeta_M(s)$ & Epstein-Zetafunktion des Gitters $M$;
  $s = -1/2$ ist der Casimir-Punkt \\
\midrule
\multicolumn{2}{l}{\textit{Symmetrien}} \\
$\text{Aut}(\Gamma)$ & Automorphismengruppe des Gitters
  (längenerhaltende Selbstabbildungen) \\
$W(F_4)$, $W(D4)$ & Weylgruppen; $|W(F_4)| = 1152$,
  $|W(D4)| = 192$ \\
$A$ & Element der Ordnung 3 in $\text{Aut}(D4)$; halbzahlig, daher
  exakt über $2A$ gerechnet \\
$\det(1-A)$ & Determinante; $|\det(1-A)| = 9$ zählt die
  Orbifold-Fixpunkte \\
$Z_3$, $C_3$ & zyklische Gruppe der Ordnung 3 (Trialität bzw.\
  Faserwirkung) \\
$S_3$, $A_5$ & symmetrische Gruppe auf 3 Elementen; alternierende
  Gruppe auf 5 Elementen (ikosaedrisch) \\
$\omega$ & dritte Einheitswurzel, $\omega = e^{2\pi i/3}$;
  Phase $+120^\circ$ \\
$\chi_0,\chi_1,\chi_2$ & die drei $Z_3$-Charaktere; indizieren die
  Sektoren des Orbifolds \\
\midrule
\multicolumn{2}{l}{\textit{Hilbertraum und Faser}} \\
$\mathcal{H}$ & Zustandsraum,
  $\mathcal{H} = L^2(\text{T}^4)\otimes\mathbb{C}^3$
  (Dok.~230/282) \\
$\mathbb{C}^3$ & Faser über jedem Punkt; trägt die
  $Z_3$-Struktur \\
$3$, $3'$ & die beiden Tripletts der Verdopplung
  $6 = 3\oplus3'$ (Dok.~285) \\
$\varphi$ & Goldener Schnitt, $\varphi = (1+\sqrt5)/2$; tritt nur
  in der $A_5$-Schicht auf \\
$\theta = 2/9$ & Phasenanteil aus Dok.~293; keine Gitterinvariante
  (Kap.~F) \\
\midrule
\multicolumn{2}{l}{\textit{FFGFT-Bezüge}} \\
$\xi = 4/30000$ & geometrische Grundkonstante der FFGFT \\
$\tilde T\cdot m = 1$ & Zeit-Masse-Relation; im Spektrum als
  Aufspaltung bei $r \ne 1$ ablesbar \\
\midrule
\multicolumn{2}{l}{\textit{Statuslabel}} \\
{}[K] & Kernableitung: gerechnet und kontrolliert \\
{}[S] & Strukturaussage: konsistent, aber nicht abgeleitet \\
Satz & mathematisch bewiesen, nicht nur numerisch geprüft \\
Grenze & explizit gebuchte Reichweitenbeschränkung \\
\end{longtable}
}
\end{center}

% ============================================================
\chapter{Die Ausgangslage: Flachheit ist nicht der Unterschied}
% ============================================================

Der Torus T$^4 = \mathbb{R}^4/\Gamma$ ist \emph{immer} flach, egal
welches Gitter $\Gamma$ ihn definiert. D4-Torus und Z$^4$-Torus
haben beide Krümmung null. Der Unterschied liegt nicht in der
Krümmung, sondern in der Gitterstruktur: Kusszahl $24$ gegen $8$,
Packungsdichte $\pi^2/16$ gegen $\pi^2/32$.

Beim Übergang in $\mathcal{H} = L^2(\text{T}^4)$ wird das Gitter zur
\textbf{Indexmenge der Fourier-Moden}: Die Basis sind ebene Wellen
$e^{ik\cdot x}$, und die zulässigen $k$ liegen auf dem dualen Gitter
$\Gamma^{*}$. In diesem Sinn wird jedes Gitter im Hilbertraum
flach -- es ist keine Kugelpackung im Raum mehr, sondern ein
Punktgitter von Quantenzahlen.

\textbf{Die Geometrie geht dabei nicht verloren, sie wandert.} Der
Laplace-Operator hat Eigenwerte $|k|^2$, und die Entartung jedes
Eigenwerts ist die Anzahl der Gittervektoren dieser Länge -- die
Thetareihe des Gitters. Die 24 kürzesten D4-Vektoren werden zur
24-fachen Entartung der ersten Anregung, und die Zerlegung
$24 = 12+12$ aus Dok.~311 (FCC-Nachbarn plus Wicklungsnachbarn)
bleibt als $k_4$-Aufspaltung exakt lesbar.

Zwei Anker halten diese Übersetzung stabil:

\begin{enumerate}
  \item \textbf{D4 ist bis auf Skalierung selbstdual} -- $D4^{*}$
        ist wieder ein D4-Gitter. Die Fourier-Transformation führt
        nicht aus der Familie heraus. (Z$^4$ ist exakt selbstdual.)
  \item \textbf{Die Faser ändert am Gitter nichts.} In
        $\mathcal{H} = L^2(\text{T}^4)\otimes\mathbb{C}^3$ hängt die
        $\mathbb{C}^3$-Faser an jedem Punkt und trägt die
        $Z_3$-Struktur; das Modengitter bleibt $\Gamma^{*}$.
\end{enumerate}

\begin{equation}
  \boxed{\;\text{Geometrie}\;\longrightarrow\;
  \text{Spektrum}\;+\;\text{Faser}\;}
\end{equation}

% ============================================================
\chapter{Das Spektrum trennt, was die Krümmung nicht trennt [K]}
% ============================================================

Vollständige Schalenzählung bis $|k|^2 = 12$:

\begin{center}
\begin{tabular}{rrrc}
\toprule
Norm $|k|^2$ & Z$^4$ & D4 & D4-Split ($k_4{=}0$ / $k_4{\ne}0$) \\
\midrule
$1$ & $8$ & $0$ & --- \\
$2$ & $24$ & $24$ & $12$ / $12$ \\
$3$ & $32$ & $0$ & --- \\
$4$ & $24$ & $24$ & $6$ / $18$ \\
$5$ & $48$ & $0$ & --- \\
$6$ & $96$ & $96$ & $24$ / $72$ \\
$7$ & $64$ & $0$ & --- \\
$8$ & $24$ & $24$ & $12$ / $12$ \\
$9$ & $104$ & $0$ & --- \\
$10$ & $144$ & $144$ & $24$ / $120$ \\
$11$ & $96$ & $0$ & --- \\
$12$ & $96$ & $96$ & $8$ / $88$ \\
\bottomrule
\end{tabular}
\end{center}

Drei Ablesungen:

\textbf{(i)~Zwei flache Tori sind spektral unterscheidbar.} Z$^4$
beginnt bei Norm~1 mit Entartung~8, D4 bei Norm~2 mit Entartung~24.
Bei D4 fehlen die ungeraden Normen \emph{vollständig} -- die
Paritätsbedingung $\sum k_i$ gerade schließt sie aus. Was die
Krümmung nicht unterscheidet, unterscheidet das Spektrum sofort.

\textbf{(ii)~Unabhängige Kontrolle bestanden.} Die D4-Reihe
$24, 24, 96, 24, 144, 96 \dots$ ist die bekannte Thetareihe
$\theta_{D4} = (\theta_3^4 + \theta_4^4)/2$; die Skriptrechnung
reproduziert sie koeffizientenweise.

\textbf{(iii)~Die $12+12$-Zerlegung ist spektral adressierbar.} Die
erste Anregung trägt exakt zwölf Moden mit $k_4 = 0$ (die räumliche
FCC-Hälfte) und zwölf mit $k_4 \ne 0$ (die Wicklungshälfte). Bei
höheren Schalen ist die Aufspaltung asymmetrisch ($6/18$ bei Norm~4,
$8/88$ bei Norm~12): \textbf{Die $12/12$-Symmetrie ist eine
Eigenschaft der ersten Schale, kein Allgemeinsatz.} Das ist
ausdrücklich zu buchen -- Dok.~311 argumentiert mit $24 = 12+12$,
und diese Aussage gilt genau dort, wo sie gemacht wird.

% ============================================================
\chapter{Anisotropie: wo die Zeit-Masse-Relation ablesbar wird [K]}
% ============================================================

Wird die vierte Richtung mit anderem Radius kompaktifiziert
($r = R_4/R_{123}$, Eigenwert $E(k) = k_1^2+k_2^2+k_3^2+(k_4/r)^2$),
spaltet die erste Schale auf:

\begin{center}
\begin{tabular}{rll}
\toprule
$r$ & $E$ (Wicklung, 12 Moden) & $E$ (FCC, 12 Moden) \\
\midrule
$1{,}0$ & $2{,}000$ & $2{,}000$ (entartet: 24) \\
$1{,}1$ & $1{,}826$ & $2{,}000$ \\
$1{,}5$ & $1{,}444$ & $2{,}000$ \\
$2{,}0$ & $1{,}250$ & $2{,}000$ \\
\bottomrule
\end{tabular}
\end{center}

Die FCC-Moden bleiben bei $E = 2$ fest, die Wicklungsmoden folgen
exakt
\begin{equation}
  E_\text{Wicklung} = 1 + \frac{1}{r^2},
  \qquad
  \Delta = 1 - \frac{1}{r^2}.
\end{equation}
Die Aufspaltung ist ein direktes Maß für das Radienverhältnis --
\textbf{die Stelle, an der die $\tilde T\cdot m = 1$-Struktur im
Spektrum ablesbar wird}: Massenrichtung und Raumrichtungen tragen
verschiedene Anregungsenergien, sobald sie verschieden lang sind.

\section{Die Kreuzungskarte ist diskret und kurz}

Alle Niveaus haben die Form $E = a + b/r^2$ mit ganzzahligen
$a, b$. Im Bereich $1 < r \le 2{,}5$ gibt es unter allen Niveaus bis
Norm~8 \textbf{exakt zwei} Kreuzungsradien:

\begin{center}
\begin{tabular}{ll}
\toprule
$r = \sqrt{2}$ & reine Wicklung $(0,0,0,\pm2)$ kreuzt die FCC-12 \\
$r = \sqrt{3}$ & reine Wicklung kreuzt die Wicklungs-12 \\
\bottomrule
\end{tabular}
\end{center}

Die Niveaufolge ändert sich also nicht kontinuierlich-chaotisch,
sondern an zwei ausgezeichneten Wurzelradien. Jede Ordnungsaussage
über die niedrigsten Anregungen zerfällt in genau drei Regime:
$r < \sqrt2$, $\sqrt2 < r < \sqrt3$, $r > \sqrt3$. Jenseits von
$\sqrt3$ ist der \emph{Grundton} des Spektrums eine reine Wicklung
(bei $r = 2$: $E = 1{,}00$ mit 2 Moden vor $1{,}25$ mit 12 vor
$2{,}00$ mit 12).

% ============================================================
\chapter{Zeit aus dem Massenkreis: die Reziprozität ist modenweise exakt [K]}
% ============================================================

Die vierte Richtung ist die Zeit-Masse-duale; ausgerollt wird auf
Zeit. Zwei Prüffragen: Wie entstehen Zeitperioden -- und gilt
$\tilde T\cdot m = 1$ nur im Mittel, wenn sich Masse und Zeit
gleich verteilen?

\section{Zeitperioden als Überlagerungsbuchhaltung}

Das Ausrollen der kompakten vierten Richtung ist die universelle
Überlagerung $\mathbb{R}\to S^1$; die Decktransformationsgruppe
ist $\mathbb{Z}$, die abgerollte Koordinate zählt
Zyklusdurchläufe. Eine $k_4$-Mode $e^{ik_4x_4/R_4}$ wird dabei
zur periodischen Funktion mit Periode $T_k = 2\pi R_4/k_4$. Da
$\pi_1(S^1) = \mathbb{Z}$ abelsch ist, ist diese Buchhaltung
\emph{verlustfrei} -- dieselbe Konstruktion, die bei
nichtabelschen Fundamentalgruppen zur Treuefrage wird, ist hier
per Struktur treu.

\section{Modenweise Exaktheit -- und was genau exakt ist}

Im Deformationsspektrum (Kap.~C) ist
$E(k) = k_1^2+k_2^2+k_3^2 + (k_4/r)^2$ wörtlich
$E^2 = p_\text{Raum}^2 + m^2$ mit $m_k = k_4/R_4$ -- \textbf{die
Wicklungszahl ist die Massenzahl}. Damit gilt
\begin{equation}
  T_k \cdot m_k \;=\; \frac{2\pi R_4}{k_4}\cdot\frac{k_4}{R_4}
  \;=\; 2\pi
  \qquad\text{exakt, für jede Mode.}
\end{equation}
Die Reziprozität gilt also \emph{nicht} nur im Mittel -- sie ist
modenweise identisch erfüllt, weil Periode und Masse aus derselben
Quantenzahl $k_4$ und demselben Radius gebaut sind. Der nächste
Abschnitt präzisiert allerdings, \emph{welche} Größe hier exakt
ist: $T_k$ ist zunächst eine Bogenlänge auf dem Kreis, also ein
Ort-Intervall. Nebenbefund:
die $12+12$-Aufspaltung (Kap.~C) ist in dieser Lesart die
Trennung masselos ($k_4 = 0$, FCC) gegen massiv (Wicklung), und
die Levelkreuzung bei $r > \sqrt3$ heißt: bei starker Anisotropie
wird eine reine Massenmode zum Grundton.

\section{Die Periode ist zunächst ein Ort, nicht eine Dauer}

Die Größe $T_k = 2\pi R_4/k_4$ ist eine \emph{Bogenlänge auf dem
Kreis}: der Abstand in der Koordinate $x_4$, nach dem sich die Mode
wiederholt. Sie ist damit ein \textbf{Ort-Intervall}, keine
verstrichene Zeit. Korrekt heißt sie $\lambda_4$, und die Relation
\begin{equation}
  \lambda_4 \cdot m \;=\; \frac{2\pi R_4}{k_4}\cdot\frac{k_4}{R_4}
  \;=\; 2\pi
\end{equation}
ist die \textbf{de-Broglie-Relation} auf der Massenrichtung --
Wellenlänge mal Wellenzahl --, nicht bereits eine Zeitaussage.
Das ist exakt und immer gültig, aber es ist eine Aussage über
$x_4$.

\textbf{Woher dann Dauer?} Erst aus der ausgerollten Koordinate.
Ein Punkt der ausgerollten Zeit ist ein Paar
\begin{equation}
  t \;=\; w\cdot 2\pi R_4 \;+\; x_4,
  \qquad w \in \mathbb{Z} \;\;(\text{Wicklung}),\;\;
  x_4 \in [0, 2\pi R_4),
\end{equation}
also \emph{Umlaufzahl} plus \emph{Ort auf dem Zyklus}. Die Dauer
steckt vollständig im Zähler $w$ -- der Decktransformation der
Überlagerung --, der Ort im Rest. Auf dem Zyklus selbst gibt es
keinen ausgezeichneten Anfang; jeder \enquote{Moment} ist eine
Stelle, die wiederkehrt (vgl.\ Dok.~313).

\section{Die zeitliche Ablesung trägt einen \texorpdfstring{$\gamma$}{gamma}-Faktor}

Fragt man nach der \emph{Schwingungsdauer in der ausgerollten
Zeit}, ist die zuständige Größe $T_\text{osc} = 2\pi/E$ mit
$E = \sqrt{p_\text{Raum}^2 + m^2}$. Dann gilt
\begin{equation}
  T_\text{osc}\cdot m \;=\; 2\pi\,\frac{m}{E}
  \;=\; \frac{2\pi}{\gamma}
  \;\le\; 2\pi,
  \qquad\text{Gleichheit} \iff p_\text{Raum} = 0 .
\end{equation}

\begin{center}
\begin{tabular}{rrrr}
\toprule
$p_\text{Raum}/m$ & $\gamma$ & $T_\text{osc}m/2\pi$ &
  $\lambda_4 m/2\pi$ \\
\midrule
$0$   & $1{,}00000$ & $1{,}0000000$ & $1{,}0000000$ \\
$0{,}5$ & $1{,}11803$ & $0{,}8944272$ & $1{,}0000000$ \\
$1$   & $1{,}41421$ & $0{,}7071068$ & $1{,}0000000$ \\
$2$   & $2{,}23607$ & $0{,}4472136$ & $1{,}0000000$ \\
$5$   & $5{,}09902$ & $0{,}1961161$ & $1{,}0000000$ \\
\bottomrule
\end{tabular}
\end{center}

\textbf{Ort und Dauer sind hier nicht dasselbe.} Exakt $2\pi$ ist
die \emph{Ort}-Ablesung; die \emph{Zeit}-Ablesung trägt den Faktor
$1/\gamma$ und fällt nur für die ruhende Mode mit ihr zusammen.
$T_k$ als \enquote{Zeitperiode} zu führen und
$T_k\cdot m = 2\pi$ als zeitliche Reziprozität zu buchen, wäre
demnach eine Verwechslung.

\textbf{Damit stehen drei Aussagen in drei Richtungen:}
\begin{center}
\begin{tabular}{lll}
\toprule
Ablesung & Relation & Gleichheit bei \\
\midrule
Ort ($\lambda_4$, de Broglie) & $=2\pi$ & immer \\
Zeit, bewegte Mode & $=2\pi/\gamma \le 2\pi$ &
  $p_\text{Raum} = 0$ \\
Massengemisch (Jensen) & $\ge 2\pi$ & scharfe Schale \\
\bottomrule
\end{tabular}
\end{center}

\noindent Die exakte Relation $\tilde T\cdot m = $ const gilt also
für die \textbf{ruhende, scharfe} Mode -- dort und nur dort
verschwinden beide Korrekturfaktoren gleichzeitig. Bewegung
drückt das Produkt herunter, Massenstreuung hebt es an. Das ist
kein Widerspruch zu Dok.~306, sondern die Kaluza-Klein-Ablesung
seiner Gültigkeitsbedingungen.

\section{Im Gemisch wird die Gleichung zur Ungleichung}

Überlagert ein Zustand mehrere $k_4$-Schalen mit Gewichten $p_k$,
gilt $\langle T\rangle\langle m\rangle =
2\pi\,\langle k\rangle\langle 1/k\rangle$ und nach Jensen
\begin{equation}
  \langle k\rangle\langle 1/k\rangle \;\ge\; 1,
  \qquad\text{Gleichheit} \iff \text{genau eine Schale besetzt.}
\end{equation}
Numerisch (Skript~4): scharfe Schale (jedes $k$) $\to$ exakt
$1{,}0000000000$; $k = 1,2$ gleichbesetzt $\to 1{,}125$;
Gleichverteilung $k = 1\ldots50 \to 2{,}295$; thermisch
$p_k \propto e^{-k/T} \to 1{,}0000227$ ($T = 0{,}1$) bis
$2{,}600$ ($T = 10$); 2000 Zufallsverteilungen: Minimum
$1{,}077$, nie unter~1.

\textbf{Buchung -- die Vermutung dreht sich um:} Die Bedingung
für $\tilde T\cdot m = 1$ ist \emph{nicht} Gleichverteilung,
sondern \textbf{Schärfe}: genau eine besetzte Massenschale.
Gleichverteilung ist im Gegenteil ein Zustand mit deutlichem
Reziprozitäts-Exzess ($2{,}29$ bei $k \le 50$). Der Exzess
$\langle k\rangle\langle 1/k\rangle - 1$ misst die Massenstreuung
des Zustands (für kleine Streuung
$\approx \mathrm{Var}(k)/\langle k\rangle^2$): ein Gemisch hat
weder scharfe Periode noch scharfe Masse, und das Produkt der
Mittel liegt strikt über $2\pi$.

\textbf{Zusammen mit dem $\gamma$-Faktor} ergibt sich die volle
Bedingung: $\tilde T\cdot m$ ist genau dann konstant, wenn die
Mode \emph{ruht} (kein $\gamma$) \emph{und} scharf ist (kein
Jensen-Exzess). Beide Abweichungen haben verschiedene Vorzeichen
und verschiedene Ursachen -- Bewegung im Raum gegen Streuung in
der Masse.

\textbf{Epistemischer Status:} gitterseitige Konsistenzprüfung der
Modenstruktur (Kaluza-Klein-Lesart der vierten Richtung); die
native FFGFT-Formulierung der Reziprozität ist Dok.~306 -- hier
wird nichts neu abgeleitet, sondern auf Verträglichkeit geprüft:
bestanden, mit der Präzisierung \enquote{exakt pro Mode, $\ge$ im
Gemisch}.

\section{Die Schließungsgabelung ist spektral unsichtbar}

Dok.~313 (nach Dok.~295) stellt die Schließungsfrage als Gabelung:
der Defekt $d = 100\,\xi = 1/75$ je Umlauf lässt drei Fälle für die
\emph{ausgerollte} Zeit zu. Mit der Zyklenlänge
$\tau_c = 2\pi/H_0 = 91{,}9$~Gyr (Dok.~312):

\begin{center}
\begin{tabular}{llll}
\toprule
Fall & Bedingung & Schließung & Periode \\
\midrule
A & $d = 0$ & nach 1 Umlauf & $91{,}9$~Gyr \\
B & $\xi$ eingefroren & nach 75 Umläufen & $6892$~Gyr \\
C & $\xi$ läuft (Rekursion) & nie & Log-Spirale \\
\bottomrule
\end{tabular}
\end{center}

\noindent Fall~B folgt aus der rationalen Rotationszahl
$1-d = 74/75$: da $\gcd(74,75) = 1$, schließt die Bahn nach genau
$75$ Umläufen, also $75\,\tau_c = 6892{,}5$~Gyr. Fall~C entsteht,
wenn $\xi$ über die Rekursion läuft; der kumulierte Defekt wird
logarithmisch und die Windung zur äquiangularen Spirale.

\textbf{Diese Tabelle wird hier nur zitiert.} Sie steht, damit die
Fälle ohne Nachschlagen verständlich sind; ihre \emph{Deutung} --
Entropiepflicht, Poincaré-Diskrepanz, thermische Geschichte --
gehört zu Dok.~312/313 und wird hier nicht geführt. Zwei Punkte
sind aber nötig, um die Zahlen nicht zu überlesen.

\textbf{Erstens: Es sind Fristen, keine Wiederholungen.} $91{,}9$
und $6892$~Gyr sagen nicht, dass sich etwas wiederholt, sondern
innerhalb welcher Frist eine Wiederkehr stattfinden \emph{müsste},
wenn der jeweilige Fall zuträfe. Geschlossen ist die
\emph{Geometrie}; ob die \emph{Geschichte} zurückkehrt, ist davon
unabhängig und nach Dok.~313 grobkörnig offen. Der
\enquote{Urknall} ist in dieser Lesart ohnehin ein \emph{Ort} auf
dem Zyklus -- die Region kleinster Massen und langsamster Uhren
--, kein Ereignis; ein geschlossener Kreis hat keinen
Anfangspunkt.

\textbf{Zweitens, und das begrenzt die Reichweite aller
Zeitangaben: Zustände an anderen Orten des Zyklus sind nicht aus
dem heutigen ableitbar.} Die übliche Rückwärtsextrapolation
($T(z) = T_0(1+z)$ aus adiabatischer Kette) setzt \emph{eine}
durchgehende thermodynamische Geschichte voraus. In der
Ortslesart ist eine frühe Epoche kein früherer Zustand desselben
Systems, sondern eine Stelle mit \emph{anderer lokaler Struktur}
-- andere Massen, andere Uhrenraten -- und damit einem eigenen
Gleichgewicht, das lokal bestimmt ist und nicht durch Fortsetzung
des heutigen. Der Korpus zeigt das an sich selbst: $T_0$ wird
nicht extrapoliert, sondern strukturell aus $\xi$ gewonnen
($k_BT_0 = (16/9)\xi$, Dok.~061/313). Genau das ist der Grund,
warum die grobkörnige Verkettung offen bleibt (Poincaré-Diskrepanz
$10^{104}$~dex, D(ii) in Dok.~313).

\noindent \textbf{Was das \emph{nicht} heißt:} Es entwertet keine
strukturelle Ableitung. Die $\Omega_m^{*}$-Kette in Dok.~313 läuft
über $\xi$ und $K_\text{frak}$, nicht über eine adiabatische
Rückrechnung, und ist von dieser Grenze nicht betroffen. Betroffen
sind allein Aussagen der Form \enquote{wie heiß war es in
Epoche~X}, sofern sie durch Fortschreiben des heutigen
Gleichgewichts gewonnen werden -- solche Aussagen bräuchten eine
eigene strukturelle Herleitung \emph{an jenem Ort}.

Naheliegend wäre, den Defekt als \emph{Twist} der Randbedingung
zu lesen ($k_4 \to k_4 + 1/75$, Scherk--Schwarz-artig). Das ist
\textbf{falsch}, und der Grund ist strukturell:

\textbf{Im eingerollten Zustand gibt es keine fraktalen
Korrekturen.} Die Eindeutigkeit der Mode auf dem Kreis erzwingt
$k_4 \in \mathbb{Z}$ -- ganzzahlig oder gar nicht. Eine
Wicklungszahl ist \emph{topologisch}; sie hat keinen Platz für
$1/75$. Korrekturen brauchen etwas, das sich \emph{akkumuliert},
und akkumulieren kann nur ein durchlaufener Weg. Schon die
Formulierung \enquote{Defekt je Umlauf} setzt das Zählen von
Umläufen voraus -- und Umläufe zählt man nur in der ausgerollten
Koordinate $w$ (Kap.~D1). Die fraktale Wegkorrektur
$K_\text{frak}$ ist eine Eigenschaft der \emph{Durchquerung},
nicht der Periodizität des Feldes.

\textbf{Daraus folgt die Zuordnung:}

\begin{center}
\begin{tabular}{lll}
\toprule
Größe & Zustand & fraktale Korrektur \\
\midrule
Thetareihe, Entartungen, $24 = 12+12$ & eingerollt & nein \\
$\text{Aut}(D4)$, Trialität, Orbifold & eingerollt & nein \\
Zirkulant, Container $4\times6$ & eingerollt & nein \\
Casimir, Wärmekern, Kongruenzen & eingerollt & nein \\
$\lambda_4\cdot m = 2\pi$ (de Broglie) & eingerollt & nein \\
\midrule
Zyklenlänge, Defekt $d$, Fall A/B/C & ausgerollt & ja \\
Umlaufzahl $w$, Geschichte & ausgerollt & ja \\
$T_\text{osc}$, $\gamma$-Faktor & ausgerollt & ja \\
\bottomrule
\end{tabular}
\end{center}

\textbf{Zwei Konsequenzen.} Erstens: \emph{Die Gabelung hat keine
spektrale Signatur.} Das Modenspektrum ist ein Objekt des
eingerollten Zustands; ob die Bahn nach einem, nach 75 oder nach
keinem Umlauf schließt, ändert an den Eigenwerten des
Laplace-Operators auf T$^4$ nichts. Die erste Schale trägt in
allen drei Fällen 24 Moden.

Zweitens, und das ist die eigentliche Entlastung: \textbf{Die
Kapitel~E--H sind fallunabhängig.} Orbifold, Zirkulant und
Container hängen nicht daran, welcher der drei Fälle zutrifft --
sie sind Aussagen über die kompakte Geometrie, und dort greift
der Defekt nicht. Die Gabelung ist eine Frage der
\emph{Geschichte} (Entropie, Wiederkehr -- Dok.~313), nicht der
\emph{Struktur}.

\textbf{Dieselbe Grenze regelt die fraktale Frage.} $D_f$ und
$K_\text{frak}$ sind \emph{verschiedene} Dinge (Dok.~133):
$D_f^\text{Raum} = 3-\xi$ ist lokal und praktisch wirkungslos
($6{,}7\times10^{-5}$), $K_\text{frak} = 1-100\xi$ ist kumulativ
über den RG-Lauf und wirkt nur dort, wo Skalen \emph{durchlaufen}
werden. Das eingerollte Spektrum trägt keines von beidem merklich
-- Kap.~J führt das aus.

\textbf{Randnotiz zur Zahl.} $75 = 1/(100\xi)$; ihre
Faktorisierung $3\cdot5^2$ trägt \emph{keine} Bedeutung -- die
Fünf hier hat nichts mit der Fünffachsymmetrie aus Kap.~F zu tun.

% ============================================================
\chapter{Umorganisationen: drei Ebenen und eine Grenze}
% ============================================================

Flache Gitter lassen sich umorganisieren -- auf drei Ebenen, die
sehr Verschiedenes kosten.

\textbf{Basiswechsel (kostenlos).} Jede unimodulare Transformation
aus $GL(4,\mathbb{Z})$ beschreibt dasselbe Punktgitter mit anderen
Erzeugern. Das Spektrum ändert sich um nichts. Reine Buchhaltung.

\textbf{Gittersymmetrien (kostenlos, strukturtragend).}
Umorganisationen, die das Gitter auf sich abbilden -- hier sitzt der
zentrale Befund, Kap.~E.

\textbf{Umbau zu einem anderen Gitter (kostet Physik).} Diskret über
Klebevektoren entlang der Kette
\begin{equation}
  D4 \;\subset\; \text{Z}^4 \;\subset\; D4^{*}
  \qquad\text{(jeweils Index 2)},
\end{equation}
oder kontinuierlich durch Deformation im Modulraum
$GL(4,\mathbb{R})/(O(4)\times GL(4,\mathbb{Z}))$. Jede Stufe
verschiebt Entartungen; D4 ist darin ein ausgezeichneter Punkt
(Kusszahl 24, in 4D nachweislich nicht überbietbar).

\textbf{Die Grenze der Spektralaussage.} Ab Dimension~4 existieren
Paare flacher Tori, die \emph{isospektral, aber nicht isometrisch}
sind (Schiemann). Der Allgemeinsatz \enquote{Spektrum bestimmt das
Gitter} ist in 4D \emph{falsch}. Für D4 selbst unkritisch -- die
Kusszahl~24 zeichnet es eindeutig aus --, aber als Regel darf es
nicht gebucht werden.

% ============================================================
\chapter{Die Trialität ist eingebaut, nicht angeheftet [K]}
% ============================================================

\section{Am Gitter nachgerechnet}

\begin{center}
\begin{tabular}{lr}
\toprule
$|\text{Aut}(D4)|$ & $1152 = |W(F_4)|$ \\
$|\text{Aut}(\text{Z}^4)|$ & $384 = 2^4\cdot4!$ \\
$|\text{Aut}(D4)|/|W(D4)|$ & $1152/192 = 6 = |S_3|$ \\
\bottomrule
\end{tabular}
\end{center}

Beide Ordnungen wurden durch Brute-Force-Zählung aller
Basis-Tupel mit identischer Gram-Matrix bestimmt und treffen die
bekannten Werte. Der Quotient $6 = |S_3|$ weist die Trialität
\emph{direkt an der Gitter-Automorphismengruppe} nach, nicht nur am
Dynkin-Diagramm: D4 ist das einzige Gitter der $D_n$-Familie mit
drei gleichberechtigten Dynkin-Armen; bei $D5, D6, \dots$ schrumpft
die äußere Symmetrie auf $Z_2$.

\section{Die 80 Ordnung-3-Elemente in drei Klassen}

Exakt gerechnet (die Trialitätselemente sind \emph{halbzahlig} --
die Klebevektor-Mischung $(\frac12,\frac12,\frac12,\frac12)$ geht
ein; Rundung auf ganze Zahlen zerstört sie, deshalb wird mit $2A$
gerechnet):

\begin{center}
\begin{tabular}{lrl}
\toprule
Klasse & Anzahl & Charakter \\
\midrule
Spur $-2$, halbzahlig & $16$ & fixpunktfrei in $\mathbb{R}^4$;
  Eigenwerte $\{\omega,\bar\omega\}$ doppelt \\
Spur $1$, ganzzahlig & $32$ & Fixebene in $\mathbb{R}^4$,
  frei auf den 24 Minimalvektoren \\
Spur $1$, halbzahlig & $32$ & 6 Fixvektoren $+$ 6 Dreierorbits \\
\bottomrule
\end{tabular}
\end{center}

\section{Zwei Konsequenzen}

\textbf{(a)~Der Orbifold ist Gittersymmetrie, keine Zusatzannahme.}
Für jedes der 16 Elemente mit Spur $-2$ gilt exakt
\begin{equation}
  |\det(1-A)| = 9,
\end{equation}
der Quotient T$^4/\langle A\rangle$ ist also ein $Z_3$-Orbifold mit
\emph{genau 9 Fixpunkten} -- der Standardzahl $3^2$. Alle neun
liegen in Drittelkoordinaten. \textbf{Die T$^4/Z_3$-Struktur muss
nicht gesetzt werden; sie ist in D4 bereits vorhanden.}

\textbf{(b)~Auf jedem Modenorbit wirkt die $Z_3$ als Zirkulant.} Die
unitäre Wirkung auf $L^2(\text{T}^4)$ permutiert die drei ebenen
Wellen eines Dreierorbits zyklisch; die Permutationsmatrix hat
Eigenwerte $\{1,\omega,\omega^2\}$, Phasen $0^\circ$, $+120^\circ$,
$-120^\circ$ -- \emph{strukturell exakt die
Zirkulant-Diagonalisierung der $\mathbb{C}^3$-Faser} (Dok.~291). Die
erste Schale liefert acht Kopien: $24 = 8\times3
= 8\times(\chi_0\oplus\chi_1\oplus\chi_2)$.

% ============================================================
\chapter{Der Negativbefund: 5 teilt 1152 nicht [K]}
% ============================================================

Die Ordnungsverteilung in $\text{Aut}(D4)$, exakt bestimmt:

\begin{center}
\begin{tabular}{lrrrrrrr}
\toprule
Ordnung & $1$ & $2$ & $3$ & $4$ & $6$ & $8$ & $12$ \\
\midrule
Anzahl & $1$ & $139$ & $80$ & $228$ & $464$ & $144$ & $96$ \\
\bottomrule
\end{tabular}
\end{center}

Summe $1152$. \textbf{Ordnung 5: null Elemente} -- und das ist kein
Rechenergebnis, sondern ein Satz: $1152 = 2^7\cdot3^2$, und $5$
teilt das nicht (Lagrange).

\textbf{Konsequenz für $\theta = 2/9$.} Die Fünffachdrehung $R_5$,
aus der $p_0 = |\langle v_0|R_5 v_e\rangle|^2 = 2/9$ folgt
(Dok.~293), ist mit der D4-Gitterstruktur \emph{prinzipiell}
unverträglich. Der Test wurde zusätzlich direkt geführt: Alle
Invarianten der linearen Gitterwirkung an den neun Orbifold-Fixpunkten
($x\!\cdot\!x$, $k\!\cdot\!x$, $x\!\cdot\!y$, sämtlich mod 1 in exakter
Bruchrechnung) haben das Nennerspektrum $\{1,2,3,6\}$ --
\textbf{keine Neuntel}. Bemerkenswert dabei: $x\!\cdot\!y$ ist zwar
ein Neuntel-Vielfaches, aber die Zähler sind stets durch 3 teilbar;
die D4-Geometrie \emph{unterdrückt} Neuntel auf der ersten Stufe.

\textbf{Was das heißt -- und was nicht.} $2/9$ ist keine Invariante
der Gitterwirkung. Das entwertet Dok.~293 nicht, sondern grenzt es
ab: Die Phase gehört einer Schicht \emph{über} dem Gitter (der
$A_5$-Struktur), das Gitter trägt die $Z_3$. Die Arbeitsteilung ist
damit exakt benannt statt vermutet.

% ============================================================
\chapter{Die Verdopplung ist dennoch vorbereitet [K]}
% ============================================================

Widerspricht der Negativbefund der Verdopplung $6 = 3\oplus3'$ aus
Dok.~285? Nein -- in drei präzisierbaren Stufen.

\textbf{(1)~Das Verbot trifft die falsche Ebene.} $5\nmid1152$
verbietet Ordnung-5-\emph{Gitterautomorphismen}. Die Verdopplung
lebt auf der \textbf{Faser} ($\mathbb{C}^3\to\mathbb{C}^6$); $A_5$
wird nirgends als Raumsymmetrie gefordert. Die Kompatibilitäts\-bedingung
ist allein, dass die Schnittstelle beider Ebenen existiert -- und
die ist exakt bestimmbar: Elementordnungen in $A_5$ sind
$\{1,2,3,5\}$, in $\text{Aut}(D4)$ sind sie
$\{1,2,3,4,6,8,12\}$; der Schnitt ist $\{1,2,3\}$, kanonisch
also $C_3$. \textbf{Genau die stellt die Trialität bereit.} Die
Schnittstelle ist nicht nur erlaubt, sie ist die einzig mögliche --
und sie ist besetzt.

\textbf{(2)~Das Gitter bereitet die Verdopplung vor.} Die zentrale
Involution $-1 \in \text{Aut}(D4)$ kommutiert mit der
Trialitäts-$Z_3$ und paart die acht Dreierorbits der ersten Schale
in \textbf{vier disjunkte Paare}. Kein Orbit ist selbstgepaart, und
das beweisbar: $-k = A^j k$ erforderte den Eigenwert $-1$ für ein
Element der Ordnung~3 -- unmöglich. Also
\begin{equation}
  24 \;=\; 8\times3 \;=\; 4\times6
  \qquad\text{(Orbit $\oplus$ Spiegelorbit).}
\end{equation}
Jeder Sechserblock trägt unter $C_3$ den Charakterinhalt
$2\times(\chi_0\oplus\chi_1\oplus\chi_2)$ -- exakt die Restriktion
von $3\oplus3'$. Die erste Anregungsschale liefert vier fertige
Container.

\textbf{(3)~Die Grenze, sauber gebucht.} Die \emph{Etikettierung} --
welcher Faktor $3$ und welcher $3'$ ist -- ist gitterblind. Die
$A_5$-Charaktere der beiden Tripletts unterscheiden sich nur auf den
5er-Klassen ($\varphi$ gegen $1-\varphi$); auf der 3er-Klasse sind
beide null, die $C_3$-Restriktionen identisch. Die Unterscheidung
hängt exakt an der Struktur, die im Gitter nicht existiert --
vollständig konsistent mit Kap.~F.

\textbf{Bilanz:} Das Gitter liefert Container und $C_3$-Inhalt, die
$A_5$-Schicht liefert die $3/3'$-Etiketten, $\varphi$ und die
Fünffach -- und damit $2/9$.

% ============================================================
\chapter{Ankerlose Ergebnisse [K]}
% ============================================================

Alle folgenden Größen sind dimensionslos -- Verhältnisse, Zählungen,
Kongruenzen. Kein Kammerton, keine SI-Skala geht ein.

\section{Zwei Kongruenzsätze für die D4-Thetareihe}

Alle vollständigen Schalen bis Norm~120 geprüft:
\begin{equation}
  N(n) \equiv 0 \pmod 3
  \qquad\text{und}\qquad
  N(n) \equiv 0 \pmod 6
  \qquad\text{ausnahmslos.}
\end{equation}
Beides ist Satz, nicht Empirie: Die Trialitäts-$Z_3$ ist auf
$\Gamma\setminus\{0\}$ fixpunktfrei (Eigenwerte $\omega,\bar\omega$
-- kein Eigenwert 1), also zerfällt \emph{jede} Schale in
Dreierorbits; zusammen mit der $-1$-Paarung ohne Selbstpaarung in
Sechserblöcke.

\section{Charakteraufgelöstes T$^4/Z_3$-Spektrum}

Folge daraus: exakte Gleichverteilung auf jeder Schale.

\begin{center}
\begin{tabular}{rrrrr}
\toprule
Norm & $N(D4)$ & $\chi_0$ & $\chi_1$ & $\chi_2$ \\
\midrule
$2$ & $24$ & $8$ & $8$ & $8$ \\
$4$ & $24$ & $8$ & $8$ & $8$ \\
$6$ & $96$ & $32$ & $32$ & $32$ \\
$10$ & $144$ & $48$ & $48$ & $48$ \\
$14$ & $192$ & $64$ & $64$ & $64$ \\
$18$ & $312$ & $104$ & $104$ & $104$ \\
\bottomrule
\end{tabular}
\end{center}

Der invariante Sektor des Orbifolds ($\chi_0$) hat exakt $N(n)/3$
Zustände pro Schale; die $Z_3$-Projektion kostet auf keiner Schale
Asymmetrie.

\textbf{Ehrliche Konsequenz.} Gleichverteilung und Sechser-Container
sind \emph{keine Auszeichnung der ersten Schale}, sondern Satz für
alle. \textbf{Die Struktur ist überall -- und kann darum nichts
selektieren.} Das ist exakt die Muster-gegen-Besetzung-Logik aus
Dok.~293: Das Muster ist skaleninvariant präsent; welche Schale
physikalisch besetzt ist, entscheidet etwas anderes.

\section{Casimir: das dichteste Gitter verliert}

$\zeta$-regularisierte Vakuumenergie eines skalaren Feldes,
$E = \pi\,\zeta_M(-1/2)$, beide Tori auf Kovolumen~1 normiert
(Epstein-Zeta über die Riemann-Split-Darstellung, dreifach
verifiziert -- gegen Direktsumme bei $s = 4, 5$ und gegen die
geschlossene Form $Z_{\text{Z}^4}(s) = 8(1-4^{1-s})\zeta(s)\zeta(s-1)$,
Übereinstimmung auf 25 Stellen):

\begin{center}
\begin{tabular}{lrr}
\toprule
 & $\zeta_M(-1/2)$ & $E = \pi\zeta_M(-1/2)$ \\
\midrule
Z$^4$-Torus & $-0{,}2966893$ & $-0{,}9320768$ \\
D4-Torus & $-0{,}2764778$ & $-0{,}8685805$ \\
\midrule
Differenz & & $+0{,}0634963$ \\
\bottomrule
\end{tabular}
\end{center}

\textbf{Der Z$^4$-Torus liegt tiefer.} Der Grund ist die
Spektrallücke: Das dichteste Gitter D4 hat bei gleichem Volumen die
\emph{größte} erste Eigenfrequenz ($\sqrt2$ gegen $1$) und damit die
\emph{weniger} negative Nullpunktsenergie. Bemerkenswert ist die
Ordnungsumkehr: Für $s > 0$ minimiert D4 die Epstein-Zeta (die
bekannte Extremalität des dichtesten Gitters), bei $s = 0$ sind alle
Gitter gleich ($\zeta(0) = -1$ universell), und bei $s = -1/2$ ist
die Ordnung invertiert.

\textbf{Interpretationsgrenze, gebucht:} Das gilt für das skalare,
periodische Feld. Anderer Feldinhalt oder andere Randbedingungen
können das Vorzeichen der Bevorzugung ändern --
\enquote{die Natur wählt Z$^4$} folgt daraus \emph{nicht}.

\section{Wärmekern: die Beobachtbarkeitsgrenze beziffert}

Für beide Tori gilt $t^2\theta(t) \to 1$ (identischer Weyl-Term bei
gleichem Volumen); der Gitterunterschied ist exponentiell klein:

\begin{center}
\begin{tabular}{rr}
\toprule
$t$ & relative Differenz \\
\midrule
$0{,}5$ & $1{,}2\times10^{-2}$ \\
$0{,}2$ & $1{,}2\times10^{-6}$ \\
$0{,}1$ & $1{,}8\times10^{-13}$ \\
$0{,}05$ & $4{,}1\times10^{-27}$ \\
\bottomrule
\end{tabular}
\end{center}

\textbf{Jede Observable, die nur auf den führenden
Wärmekern-Koeffizienten beruht, ist gitterblind.} Der
D4/Z$^4$-Unterschied sitzt ausschließlich in der exponentiell
unterdrückten Thetareihen-Feinstruktur -- beziehungsweise direkt in
der Spektrallücke.

% ============================================================
\chapter{Was hier noch nackt ist -- und was nicht [K]}
% ============================================================

Dieses Dokument rechnet mit dem \emph{glatten} Laplace-Operator auf
$\text{T}^4 = \mathbb{R}^4/\Gamma$. Zu klären ist, was die FFGFT
an fraktaler Struktur dagegenhält -- und dabei sind \textbf{zwei
Größen sauber zu trennen}, die leicht verwechselt werden:

\begin{center}
\begin{tabular}{lll}
\toprule
Größe & Charakter & Wirkung \\
\midrule
$D_f^\text{Raum} = 3-\xi$ (Dok.~133) & \emph{lokal} &
  $6{,}67\times10^{-5}$ \\
$K_\text{frak} = 1-100\xi$ (Dok.~133, A040) & \emph{kumulativ} &
  $1{,}33\times10^{-2}$ \\
\bottomrule
\end{tabular}
\end{center}

\noindent Dok.~133 leitet $K_\text{frak}$ ausdrücklich \emph{nicht}
aus der lokalen Dimension her, sondern aus der \enquote{kumulativen
Wirkung} über den RG-Lauf von der Planck- zur elektroschwachen
Skala (17 Größenordnungen); der Faktor~100 ist RG-Verstärkung
(Dok.~190). Die lokale Dimension allein liefert nach Dok.~133
\enquote{eine verschwindend kleine Korrektur -- praktisch Eins},
nachgerechnet $(D/3)^{D/2} = 0{,}99993$. Das Verhältnis der beiden
Effekte ist exakt $200$.

\textbf{Zur Notation, weil hier leicht verrutscht wird:} Die
fraktale Dimension der FFGFT ist $D_f^\text{Raum} = 3-\xi$ -- die
\emph{räumliche}, nicht eine vierdimensionale. Auch der Exponent
$3/2$ in $K_\text{frak}^{3/2}$ (Dok.~018) wird als halbe
\emph{Raum}dimension begründet. Es sind drei Dimensionen, nicht
vier; eine Größe $D_f = 4-\xi$ kennt der Korpus nicht.

\textbf{Daraus folgt die Zuordnung} -- und sie deckt sich mit der
Grenze aus Kap.~D2 (eingerollt gegen ausgerollt): Ein
Laplace-Operator auf dem kompakten T$^4$ ist \emph{lokal}. Eine
kumulative Größe kann er nicht tragen; $K_\text{frak}$ kommt nicht
über den Operator herein, sondern über die \textbf{Brücken\-kon\-stan\-ten}
($E_0$, $v$) bei der Skalenumrechnung -- genau so, wie R72 es am
Beispiel der beiden $E_0$-Werte zeigt (Verhältnis $1/\sqrt{K_\text{frak}}$).
Die Konsequenz zerfällt in drei sehr ungleiche Teile.

\section{Unberührt: alles Ankerlose}

Nach der Zuordnungsregel (Register~R70, Dok.~313) greift
$K_\text{frak}$ bei absoluten Größen und kürzt sich bei
Verhältnissen derselben Ebene. Damit bleibt unberührt:

\begin{itemize}
  \item alle \textbf{Zählungen} -- $24$, $8$, $N(n)$, $8\times3$,
        $4\times6$: ganzzahlige Kombinatorik, keine Länge;
  \item die \textbf{Kongruenzen} $N(n)\equiv0 \bmod 3$ und
        $\bmod\,6$: Gruppentheorie, kein metrischer Inhalt;
  \item die \textbf{Charakteraufteilung} $N/3$ je Sektor;
  \item alle \textbf{Verhältnisse} $E/E_\text{min}$ und die
        Kreuzungsradien $\sqrt2$, $\sqrt3$ -- Struktur durch
        Struktur, der Faktor kürzt sich heraus.
\end{itemize}

\noindent Das ist der eigentliche Grund, warum Kap.~H
\enquote{ankerlose Ergebnisse} heißt: Diese Aussagen sind
fraktal-invariant.

\section{Ein Faktor: alle absoluten Energien}

Mit einem Anker wird $E = \hbar H_0\,|k|$ (Grundmode $= \hbar H_0$,
siehe Kap.~B in Verbindung mit Dok.~312). Diese Werte sind nackt und
tragen den Faktor:

\begin{center}
\begin{tabular}{lrr}
\toprule
 & nackt & mit $K_\text{frak}$ \\
\midrule
erste D4-Schale ($|k| = \sqrt2$) & $3{,}230\times10^{-52}$~J &
  $3{,}187\times10^{-52}$~J \\
reine Wicklung ($|k| = 2$) & $4{,}567\times10^{-52}$~J &
  $4{,}507\times10^{-52}$~J \\
\bottomrule
\end{tabular}
\end{center}

\noindent Durchgehend $-1{,}33\,\%$. Dasselbe gilt für die
Casimir-Energien aus Kap.~H und für die Spektrallücke selbst.
Buchhalterisch erledigt, nicht strukturell interessant.

\textbf{Wichtig für die Zuordnung:} Der Faktor sitzt hier
\emph{nicht} im Spektrum, sondern im \emph{Anker}. Die
dimensionslosen Eigenwerte $|k|^2$ bleiben unverändert; korrigiert
wird die Umrechnung in SI über $\hbar H_0$ -- eine
Skalenbrücke, also genau der Weg, auf dem eine kumulative Größe
wirken darf. Wer nur Verhältnisse betrachtet, sieht nichts davon
(A040-Regel, R72).

\section{Der Operator, nicht die Umrechnung: die Störungsrechnung [K]}

Der dritte Teil ist nicht durch Multiplikation zu erledigen. Bei
$D_f^\text{Raum} = 3-\xi$ ist der Operator \emph{kein exakt glatter
Laplace mehr} -- und exakte Symmetrie ist genau das, was die
Entartungen trägt. Maßgeblich ist hier allein der \emph{lokale}
Effekt: $K_\text{frak}$ ist nach Dok.~133 kumulativ und gehört
darum nicht auf einen lokalen Operator, $O(100\xi)$ scheidet als
Störungsstärke aus.

\begin{center}
\begin{tabular}{lrr}
\toprule
Störungsstärke & $\Delta E/E$ & $\Delta E$ bei $E = \hbar H_0$ \\
\midrule
lokal, $(D/3)^{D/2}$ (Dok.~133) & $6{,}7\times10^{-5}$ &
  $1{,}5\times10^{-56}$~J \\
$O(\xi)$ (obere Schranke) & $1{,}3\times10^{-4}$ &
  $3{,}0\times10^{-56}$~J \\
\bottomrule
\end{tabular}
\end{center}

\noindent Die Störung liegt damit \textbf{zwei Größenordnungen
unter} dem, was eine Abschätzung über $K_\text{frak}$ ergäbe.

\noindent \textbf{Die 24-fache Entartung wäre dann keine Entartung
mehr, sondern ein Multiplett mit Feinstruktur.} In welche
Untermultipletts es zerfällt, entscheidet die \emph{Symmetrie der
Störung} -- und zwar, wie die Rechnung zeigt, über die
\emph{Darstellungstheorie}, nicht über die Orbitstruktur:

\begin{center}
\begin{tabular}{ll}
\toprule
Störung respektiert & Multiplettstruktur der 24 (gerechnet) \\
\midrule
volle $\text{Aut}(D4)$ ($c_g$ nur normabhängig) &
  $9 + 8 + 4 + 2 + 1$ \\
Trialität ($Z_3$) und $-1$ & $8\times1 \;+\; 8\times2$ \\
nur $Z_3$ & $8\times1 \;+\; 8\times2$ \\
nur $-1$ & $24\times1$ \\
nichts davon & $24\times1$ \\
\bottomrule
\end{tabular}
\end{center}

\textbf{Gerechnet} (Skript 4): Störungstheorie erster Ordnung auf
der $24\times24$-Matrix $\langle k|V|k'\rangle = c_{k-k'}$,
Symmetrieklassen durch Orbit-Mittelung der Koeffizienten, je drei
unabhängige Zufallsstörer pro Klasse -- die Multiplettstruktur ist
in allen Fällen seedunabhängig identisch.

\textbf{Die naheliegende Erwartung ist falsch.} Man vermutet, die
Orbits blieben zusammen ($Z_3$ und $-1$: die 24 als Ganzes; nur
$Z_3$: 8~Dreierorbits; nur $-1$: 4~Sechserblöcke). Tatsächlich
folgen die Multiplettgrößen den \emph{irreduziblen Darstellungen}
der respektierten Symmetriegruppe, nicht den Bahnen. Konkret: (i)~Jede ortsabhängige Störung bricht die
Translationsinvarianz des Torus -- und \emph{die} trug die
24er-Entartung, nicht die Punktsymmetrie. (ii)~$Z_3$ ist abelsch,
alle Irreps eindimensional; die Dubletts sind die antiunitäre
Paarung der konjugierten Charaktere $\chi_1\leftrightarrow\chi_2$,
der $\chi_0$-Sektor zerfällt in 8~Singuletts. (iii)~$-1$ allein
erzwingt gar nichts (freie $Z_2$-Wirkung, eindimensionale
Charaktere). (iv)~Selbst der maximal symmetrische Störer hält die
24 nicht zusammen: $24 = 9\oplus8\oplus4\oplus2\oplus1$,
konsistent mit den Irrep-Dimensionen von $W(F_4)$ auf der
Wurzelschale [K numerisch; die Irrep-Zuordnung ist B, ohne
Charakterrechnung].

\textbf{Damit ist der offene Punkt geschlossen -- und entschärft:}
Qualitativ gilt weiterhin, dass von der 24er-Entartung nur so viel
stabil ist, wie die Störungssymmetrie an Irrep-Dimension hergibt
(im besten Fall $9+8+4+2+1$, generisch gar nichts). Quantitativ
liegt die zuständige Störung aber bei $\sim10^{-5}$ relativ, nicht
bei $10^{-2}$: Die Aufspaltung wäre eine Feinstruktur weit
unterhalb jeder Ablesbarkeit, und die Container-Struktur
($4\times6$, Kap.~G) bleibt praktisch intakt.

\textbf{Bilanz des Kapitels.} \emph{Nackt} sind nur die
SI-Absolutwerte, und zwar über den Anker, nicht über den
Operator. Die eingerollten Größen dieses Dokuments --
Entartungen, Verhältnisse, Kongruenzen, Container, Casimir- und
Wärmekern-Verhältnisse -- sind \textbf{nicht} nackt, sondern
korrekt: $K_\text{frak}$ greift dort strukturell nicht.
Das ist dieselbe Grenze wie in Kap.~D2, nur von der anderen
Seite gelesen.

% ============================================================
\chapter{Statusübersicht}
% ============================================================

{\small
\begin{longtable}{lp{6.4cm}l}
\toprule
Aussage & Inhalt & Status \\
\midrule
\endfirsthead
\toprule
Aussage & Inhalt & Status \\
\midrule
\endhead
\bottomrule
\endfoot
Übersetzung & Geometrie $\to$ Spektrum $+$ Faser; Gitter wird
  Modenindex, Entartungen $=$ Thetareihe & [K] \\
Unterscheidbarkeit & Z$^4$ ab Norm 1 (8), D4 ab Norm 2 (24);
  ungerade Normen bei D4 leer & [K] \\
$12+12$ & erste Schale exakt FCC $+$ Wicklung; höhere Schalen
  asymmetrisch ($6/18$, $8/88$) & [K] \\
Anisotropie & $E_\text{Wick} = 1 + 1/r^2$; Kreuzungen nur bei
  $r = \sqrt2$ und $\sqrt3$; drei Regime & [K] \\
Trialität & $|\text{Aut}(D4)| = 1152$, Quotient $6 = |S_3|$;
  80 Ordnung-3-Elemente in drei Klassen & [K] \\
Orbifold & 16 Elemente mit $|\det(1-A)| = 9$: T$^4/Z_3$ ist
  Gittersymmetrie, keine Zusatzannahme & [K] \\
Zirkulant & $Z_3$ wirkt auf jedem Dreierorbit mit
  $\{1,\omega,\omega^2\}$ $=$ $\mathbb{C}^3$-Faserstruktur & [K] \\
$5 \nmid 1152$ & keine Ordnung-5-Automorphismen; $2/9$ ist keine
  Gitterinvariante (Nenner $\{1,2,3,6\}$) & [K] Satz \\
Verdopplung & Schnittstelle $A_5 \cap \text{Aut}(D4) = C_3$,
  besetzt; $24 = 4\times6$ als Container & [K] \\
Etikettierung & $3$ gegen $3'$ ist gitterblind (Unterschied nur auf
  5er-Klassen) & [K] Grenze \\
Kongruenzen & $N(n) \equiv 0 \bmod 3$ und $\bmod 6$ auf allen
  60 geprüften Schalen, mit Beweisgrund & [K] Satz \\
Selektion & Container und Gleichverteilung sind generisch --
  die Struktur kann keine Schale auszeichnen & [K] \\
Casimir & $E(\text{Z}^4) = -0{,}9321$ gegen
  $E(D4) = -0{,}8686$; Ordnungsumkehr gegenüber $s>0$
  & [K], Feldinhalt-abhängig \\
Wärmekern & Weyl-Term identisch, Differenz $4\times10^{-27}$ bei
  $t = 0{,}05$: führende Ordnung ist gitterblind & [K] \\
Spektralsatz & \enquote{Spektrum bestimmt Gitter} in 4D falsch
  (Schiemann); für D4 unkritisch, aber keine Regel & [K] Grenze \\
Reziprozität & $\lambda_4\cdot m = 2\pi$ exakt als Ort-Relation
  (de Broglie); zeitlich $2\pi/\gamma \le 2\pi$; im Gemisch
  $\ge 2\pi$ (Jensen) -- exakt nur ruhend und scharf & [K] \\
Störungsmultipletts & Aufspaltung folgt Irrep-Dimensionen, nicht
  Orbits: $9{+}8{+}4{+}2{+}1$ (voll), $8{\times}1{+}8{\times}2$
  ($Z_3$), $24{\times}1$ (generisch) & [K] \\
Schließungsfall & eingerollt keine fraktale Korrektur
  ($k_4$ topologisch); Gabelung A/B/C spektral unsichtbar,
  Kap.~E--H fallunabhängig & [K] \\
Fraktal-Status & $D_f^\text{Raum} = 3-\xi$ lokal ($6{,}7{\times}10^{-5}$)
  gegen $K_\text{frak}$ kumulativ ($1{,}33\,\%$, RG-Lauf, Dok.~133):
  nackt sind nur SI-Absolutwerte über den Anker, nicht das
  eingerollte Spektrum & [K] \\
\end{longtable}
}

\medskip
\noindent\textbf{Bilanz:} Im Hilbertraum wird das Gitter flach, aber
nicht folgenlos. Die Kusszahl wird zur Entartung, die Trialität zum
Orbifold und zur Faser, das Radienverhältnis zur Niveauaufspaltung.
Ebenso scharf ist die \emph{Arbeitsteilung} bestimmt: Die
Fünffachsymmetrie -- und mit ihr $\varphi$ und $2/9$ -- kann das
Gitter \emph{selbst} nicht tragen, und zwar per Satz
($5 \nmid 1152$), nicht aus Mangel an Rechenaufwand. Das ist eine
Aussage über die \emph{Ebene}, nicht über die Größe: Auf der
Faser ist $2/9$ ausdrücklich zulässig, denn die Schnittstelle
$A_5 \cap \text{Aut}(D4) = C_3$ ist durch die Trialität besetzt
(Kap.~G), und das Gitter stellt mit $24 = 4\times6$ sogar die
Container für die Verdopplung $3\oplus3'$ bereit. Ausgeschlossen
ist also nicht $2/9$, sondern nur seine Herleitung \emph{aus}
Gitterinvarianten -- die Phase kommt nachweislich eine Schicht
höher herein. Und die Container-Struktur ist generisch: Sie stellt
bereit, sie selektiert nicht.

% ============================================================
\chapter{Prüfskripte}
% ============================================================

Drei Skripte, alle deterministisch, alle Kontrollen bestanden:

\begin{itemize}
  \item \texttt{d4\_skript\_1\_spektrum\_deformation.py} --
        Thetareihen D4/Z$^4$, $k_4$-Split, Deformation $r$, exakte
        Kreuzungsradien (nur Standardbibliothek).
  \item \texttt{d4\_skript\_2\_trialitaet\_orbifold\_phasen.py} --
        $\text{Aut}(D4) = 1152$, Ordnung-3-Klassifikation (mit $2A$
        exakt gerechnet), $|\det(1-A)| = 9$, Zirkulant, Phasentest
        $2/9$, Ordnungsverteilung, $24 = 4\times6$ (numpy).
  \item \texttt{d4\_skript\_3\_schalen\_casimir.py} --
        Schalensätze, Epstein-Zeta bei $s = -1/2$, Wärmekern
        (numpy, mpmath).
  \item \texttt{d4\_skript\_4\_stoerung\_reziprozitaet.py} --
        Störungstheorie 1.~Ordnung auf der 24er-Schale
        (Multiplettstruktur nach Symmetrieklasse, je 3~Seeds) und
        Reziprozität $T\cdot m$ (modenweise; Jensen im Gemisch)
        (numpy, mpmath).
\end{itemize}

Kontrollen gegen bekannte Werte: $\theta_{D4}$,
$|W(F_4)| = 1152$, $2^4\cdot4! = 384$, $|W(D4)| = 192$,
Orbifold-Fixpunktzahl $3^2 = 9$, Jacobi-Identität für
$Z_{\text{Z}^4}$.

\medskip
\noindent\textbf{Registrierung:} Aufnahme in Dok.~190 erfolgt
separat, sofern die Abgrenzung zu Dok.~293 dort gebucht werden soll.

\end{document}
